%% RSI Appendix — To be included via %% RSI Appendix — To be included via %% RSI Appendix — To be included via %% RSI Appendix — To be included via \input{appendix}
\section{Appendix: Supplementary Material}
\label{sec:appendix}

%% ─────────────────────────────────────────────────────────────
%% A.1 — HYPERPARAMETERS
%% ─────────────────────────────────────────────────────────────

\subsection{Hyperparameters}

Table~\ref{tab:hyperparams} lists all hyperparameters used in the RSI
Togolese fiscal instantiation. All values are derived from institutional
knowledge (OTR fiscal rules, historical compliance estimates) or set to
standard non-informative values. No hyperparameter is learned from
labeled data.

\begin{table}[H]
\centering
\caption{RSI hyperparameters. Prior parameters encode institutional
         knowledge; inference parameters are standard.}
\label{tab:hyperparams}
\renewcommand{\arraystretch}{1.15}
\begin{tabular}{llll}
\toprule
\textbf{Parameter} & \textbf{Value} & \textbf{Scope} & \textbf{Justification} \\
\midrule
\multicolumn{4}{l}{\textit{Prior hyperparameters (per rule)}} \\
$\alpha_i, \beta_i$ & See Table~\ref{tab:rule_priors} & Per rule & OTR historical rates \\
$\sigma_i$ (drift) & 0.5--2.0 & Per rule & Expected parametric stability \\
\midrule
\multicolumn{4}{l}{\textit{Inference parameters}} \\
$\sigma_{\mathrm{obs}}$ & 0.0625 & Global & Observation noise scale \\
$\varepsilon$ (ELBO tol.) & $10^{-8}$ & Global & Standard convergence criterion \\
Max iterations & 100 & Global & Upper bound (converges in 2) \\
$\tau$ (discretization) & 0.5 & Global & Neutral threshold \\
\bottomrule
\end{tabular}
\end{table}

\begin{table}[H]
\centering
\caption{Per-rule prior parameters.}
\label{tab:rule_priors}
\renewcommand{\arraystretch}{1.15}
\begin{tabular}{llccl}
\toprule
\textbf{Rule} & \textbf{Description} & $\alpha_i$ & $\beta_i$
  & $\mathbb{E}[c_i]$ \\
\midrule
R1\_TVA  & Value Added Tax          & 7 & 3 & 0.70 \\
R2\_IS   & Corporate Income Tax     & 6 & 4 & 0.60 \\
R3\_IMF  & Minimum Tax              & 7 & 3 & 0.70 \\
R4\_TPU  & Informal Sector Tax      & 5 & 5 & 0.50 \\
R5\_IRPP & Income Tax Withholding   & 6 & 4 & 0.60 \\
R6\_PAT  & Business License         & 7 & 3 & 0.70 \\
R7\_DECL & Annual Declaration       & 6 & 4 & 0.60 \\
R8\_BANK & Bank Formalization       & 5 & 5 & 0.50 \\
\bottomrule
\end{tabular}
\end{table}

%% ─────────────────────────────────────────────────────────────
%% A.2 — DATASET DOCUMENTATION
%% ─────────────────────────────────────────────────────────────

\subsection{RSI-Togo-Fiscal-Synthetic v2.0}

\paragraph{Overview.}
The benchmark contains 2{,}000 synthetic enterprises across two
regulatory periods (2022--2024 and 2025), encoding 8 fiscal rules.
Each enterprise record is structured in four layers: observable
features (segment, declared turnover), latent ground truth
($c_i^*$ per rule), noisy compliance signals with 18\% missing data,
and binary labels for baseline evaluation only.

\paragraph{Market segments.}
Enterprises are assigned to four segments reflecting the structure
of the Togolese market. Turnover follows a log-normal distribution
within each segment, whose aggregate approximates a power
law~\citep{gabaix2009power}.

\begin{center}
\renewcommand{\arraystretch}{1.15}
\begin{tabular}{lcccc}
\toprule
\textbf{Segment} & \textbf{Share} & \textbf{$\mu_{\log}$}
  & \textbf{$\sigma_{\log}$} & \textbf{Median CA (FCFA)} \\
\midrule
Informal      & 45\% & 16.3 & 0.7 & $\sim$12M \\
Small formal  & 25\% & 17.6 & 0.5 & $\sim$44M \\
Medium        & 18\% & 18.6 & 0.4 & $\sim$119M \\
Large         & 12\% & 20.0 & 0.5 & $\sim$485M \\
\bottomrule
\end{tabular}
\end{center}

\paragraph{Under-declaration model.}
Observed turnover is systematically lower than true turnover:
\[
  \hat{x} = x \cdot \underbrace{\beta}_{\text{systematic}} \cdot
  \underbrace{\varepsilon}_{\text{noise}}, \quad
  \beta \sim \operatorname{Beta}(7, 3), \quad
  \varepsilon \sim \mathcal{N}(1,\, 0.04^2)
\]
The Beta(7,3) distribution has mean $0.70$ and concentrates mass
between 0.50 and 0.90, consistent with empirical estimates from
Sub-Saharan African fiscal studies.

\paragraph{Compliance signals.}
For each entity $j$ and applicable rule $r_i$, the continuous
compliance signal is drawn from
$\operatorname{Beta}(c_{ji} \cdot \kappa,\, (1{-}c_{ji}) \cdot \kappa)$
with $\kappa = 10$, ensuring an unbiased signal centered on the true
entity compliance rate: $\mathbb{E}[s_{ji}^{\mathrm{raw}}] = c_{ji}$.
For R4\_TPU (informal sector), the signal is binary:
$s_{ji} \sim \operatorname{Bernoulli}(c_{ji})$.

\paragraph{Applicability and rules.}
Rules R7\_DECL and R8\_BANK apply universally. The remaining 6 rules
have turnover-based thresholds, ensuring every entity is subject to
at least 3 rules (K\,$\geq$\,3). The distribution is:
K\,=\,3 (56\%), K\,=\,5 (15\%), K\,=\,6 (10\%), K\,=\,7 (19\%).

\begin{center}
\renewcommand{\arraystretch}{1.15}
\begin{tabular}{llll}
\toprule
\textbf{Rule} & \textbf{Description} & \textbf{Condition} & \textbf{Threshold} \\
\midrule
R1\_TVA  & Value Added Tax       & CA $\geq$ threshold & 60M (p1), 100M (p2) \\
R2\_IS   & Corporate Income Tax  & CA $\geq$ threshold & 100M \\
R3\_IMF  & Minimum Tax           & CA $\geq$ threshold & 60M \\
R4\_TPU  & Informal Sector Tax   & CA $<$ threshold    & 60M \\
R5\_IRPP & Income Tax Withholding & CA $\geq$ threshold & 30M \\
R6\_PAT  & Business License      & CA $\geq$ threshold & 30M \\
R7\_DECL & Annual Declaration    & Always              & (universal) \\
R8\_BANK & Bank Formalization    & Always              & (universal) \\
\bottomrule
\end{tabular}
\end{center}

\paragraph{Regulatory change event.}
Period~1 (2022--2024) uses a VAT threshold of 60M~FCFA. Period~2
(2025) uses 100M~FCFA, reflecting Law n\textdegree{}2024-007
(30~December~2024). This provides a clean natural experiment for
validating Theorem~\ref{thm:o1}.

\paragraph{Column structure.}
Each enterprise record contains 48 columns: 6 identifiers and
features (\texttt{entity\_id}, \texttt{period}, \texttt{segment},
\texttt{true\_ca}, \texttt{obs\_ca}, \texttt{under\_decl\_ratio}),
5 columns per rule (applicability, true compliance, raw signal,
discretized signal, label), and 2 aggregate columns
(\texttt{label\_any\_nc}, \texttt{n\_applicable}).

%% ─────────────────────────────────────────────────────────────
%% A.3 — BASELINES
%% ─────────────────────────────────────────────────────────────

\subsection{Baseline Configurations}

\begin{table}[H]
\centering
\caption{Baseline model configurations.}
\label{tab:baselines}
\renewcommand{\arraystretch}{1.15}
\begin{tabular}{lll}
\toprule
\textbf{Model} & \textbf{Parameter} & \textbf{Value} \\
\midrule
LR      & solver & lbfgs \\
        & max\_iter & 1000 \\
        & feature scaling & StandardScaler \\
        & train/test split & 70\%/30\%, stratified \\[4pt]
XGBoost & n\_estimators & 150 \\
        & max\_depth & 4 \\
        & learning\_rate & 0.1 (default) \\
        & train/test split & 70\%/30\%, stratified \\[4pt]
MLP     & hidden\_layer\_sizes & (128, 64, 32) \\
        & activation & ReLU \\
        & max\_iter & 1000 \\
        & early\_stopping & True \\
        & learning\_rate & adaptive \\
        & feature scaling & StandardScaler \\[4pt]
RBS     & Logic & Flag if $s_{ji}^{\mathrm{disc}} = 0$ for any applicable rule \\
        & Missing handling & Ignore (no flag) \\
\bottomrule
\end{tabular}
\end{table}

%% ─────────────────────────────────────────────────────────────
%% A.4 — REPRODUCIBILITY
%% ─────────────────────────────────────────────────────────────

\subsection{Reproducibility}

All experiments use a fixed random seed (42) for full deterministic
reproducibility. The implementation uses Python~3.11 with NumPy,
SciPy, scikit-learn, and pandas. No GPU is required. The full
experiment pipeline completes in under 30~seconds on standard
hardware. The dataset, framework, and all experiment code are
released publicly under the MIT license.

%% ─────────────────────────────────────────────────────────────
%% A.5 — DETAILED PROOFS
%% ─────────────────────────────────────────────────────────────

\subsection{Detailed Proof of Theorem~\ref{thm:o1} ($O(1)$ Adaptability)}

\begin{proof}
Let $\mathcal{U}_k$ be a regulatory update that changes the parameters
of rule $r_k$ from $\Theta_k$ to $\Theta_k'$.

\textit{Step 1: Prior factorization.}
The RSI prior factorizes as
$P(\mathbf{S}) = \prod_{i=1}^n P(c_i \mid \Theta_i) \cdot P(\delta_i \mid \Theta_i)$.
The posterior satisfies:
\[
  P(\mathbf{S} \mid \mathbf{D}) \propto
  \prod_{i=1}^n P(\mathbf{D}_i \mid c_i) \cdot P(c_i \mid \Theta_i) \cdot P(\delta_i \mid \Theta_i)
\]
After the update $\Theta_k \to \Theta_k'$, only the $k$-th factor
changes:
\[
  P'(\mathbf{S} \mid \mathbf{D}) \propto
  P(\mathbf{S} \mid \mathbf{D}) \cdot
  \frac{P(c_k \mid \Theta_k') \cdot P(\delta_k \mid \Theta_k')}
       {P(c_k \mid \Theta_k) \cdot P(\delta_k \mid \Theta_k)}
\]

\textit{Step 2: Drift correction.}
For a threshold change $\theta_k \to \theta_k'$, define
$\varepsilon = (\theta_k' - \theta_k)/\theta_k$. The drift posterior
mean updates as $\bar{\mu}_k \leftarrow \bar{\mu}_k - \varepsilon$.
This is a single scalar subtraction: $O(1)$.

\textit{Step 3: Invariance of compliance parameters.}
The compliance rate $c_k$ is independent of the threshold magnitude
(it measures the \emph{proportion} of applicable entities that comply,
not the threshold itself). Therefore $\bar{\alpha}_k$ and
$\bar{\beta}_k$ are invariant under $\mathcal{U}_k$. Similarly,
$\bar{\sigma}_k$ depends on the number of observations and prior
variance, neither of which changes under $\mathcal{U}_k$.

\textit{Step 4: Normalization.}
Under conjugate families, the posterior normalization constant
$B(\bar{\alpha}_k, \bar{\beta}_k)$ is computed analytically and does
not require re-processing the data. The updated posterior is a
properly normalized Beta distribution.

\textit{Total cost:} one scalar subtraction ($\bar{\mu}_k$), zero
data access. $\mathcal{C}^{\mathrm{RSI}}(\mathcal{U}_k) = O(1)$
with respect to $|\mathbf{D}|$ and $n$. \qed
\end{proof}

\subsection{Detailed Proof of Theorem~\ref{thm:bvm} (BvM Consistency)}

\begin{proof}
We prove the result for a single rule $r_i$ (the extension to
$\mathbf{S}$ follows by factorization).

Let $s_1, s_2, \ldots, s_m$ be i.i.d.\ Bernoulli($c_i^*$) compliance
signals, where $c_i^* \in (0,1)$ is the true population compliance rate.

\textit{Step 1: Posterior computation.}
The conjugate posterior is
$Q^*(c_i) = \operatorname{Beta}(\bar{\alpha}_i, \bar{\beta}_i)$ with
$\bar{\alpha}_i = \alpha_i + \sum_{j=1}^m s_j$ and
$\bar{\beta}_i = \beta_i + \sum_{j=1}^m (1-s_j)$.

\textit{Step 2: Posterior mean convergence.}
Define $\bar{s} = m^{-1}\sum_j s_j$. Then:
\[
  \mathbb{E}_Q[c_i] = \frac{\bar{\alpha}_i}{\bar{\alpha}_i + \bar{\beta}_i}
  = \frac{\alpha_i + m\bar{s}}{\alpha_i + \beta_i + m}
  \to \bar{s} \to c_i^* \quad (m \to \infty)
\]
where the first convergence follows from $m \gg \alpha_i + \beta_i$
and the second from the strong law of large numbers.

\textit{Step 3: Posterior variance.}
\[
  \mathrm{Var}_Q[c_i] = \frac{\bar{\alpha}_i \bar{\beta}_i}
  {(\bar{\alpha}_i + \bar{\beta}_i)^2 (\bar{\alpha}_i + \bar{\beta}_i + 1)}
  \sim \frac{c_i^*(1-c_i^*)}{m} \quad (m \to \infty)
\]
This is the Bernoulli variance divided by $m$, confirming
$\sigma[c_i] \sim 1/\sqrt{m}$.

\textit{Step 4: Asymptotic normality.}
By the central limit theorem, $\sqrt{m}(\bar{s} - c_i^*)
\xrightarrow{d} \mathcal{N}(0, c_i^*(1-c_i^*))$. Since
$\mathbb{E}_Q[c_i] \to \bar{s}$ and $\mathrm{Var}_Q[c_i] \to
c_i^*(1-c_i^*)/m$, we obtain:
\[
  \sqrt{m}(\mathbb{E}_Q[c_i] - c_i^*) \xrightarrow{d}
  \mathcal{N}(0, \mathcal{I}(c_i^*)^{-1})
\]
where $\mathcal{I}(c_i^*) = [c_i^*(1-c_i^*)]^{-1}$ is the Fisher
information of the Bernoulli model. This is the classical
Bernstein-von Mises theorem applied to the Beta-Bernoulli conjugate
pair \citep{van2000asymptotic}. \qed
\end{proof}

%% ─────────────────────────────────────────────────────────────
%% A.6 — SIGNAL DISCRETIZATION BIAS
%% ─────────────────────────────────────────────────────────────

\subsection{Signal Discretization and Calibration Bias}

The compliance signal $s_{ji}^{\mathrm{raw}} \in [0,1]$ is unbiased
by construction:
$\mathbb{E}[s_{ji}^{\mathrm{raw}}] = c_{ji}$.
Discretization at threshold $\tau = 0.5$ introduces a systematic bias.
For an entity with true compliance $c_{ji}$, the probability that the
discretized signal equals 1 is:
\[
  P(s_{ji}^{\mathrm{disc}} = 1) =
  1 - I_{0.5}(c_{ji} \cdot \kappa,\, (1{-}c_{ji}) \cdot \kappa)
\]
where $I_x(a,b)$ is the regularized incomplete Beta function. When
$c_{ji} > 0.5$, this probability exceeds $c_{ji}$ (the discrete signal
overestimates compliance). When $c_{ji} < 0.5$, it underestimates.
The bias is largest for $c_{ji}$ far from 0.5 and decreases as
$\kappa \to \infty$.

This explains the calibration misses in Table~\ref{tab:calibration}:
rules with true population compliance far from 0.5 (R3\_IMF at 0.876,
R7\_DECL at 0.666) exhibit the largest bias, while rules near 0.5
(R4\_TPU at 0.345, R8\_BANK at 0.491) are well-calibrated.

A natural mitigation is to use continuous signals directly with a
Beta-likelihood formulation, eliminating discretization entirely while
preserving conjugacy under moment-matching approximations. This
extension is developed in future work.

%% ─────────────────────────────────────────────────────────────
%% A.7 — PRIOR SENSITIVITY ANALYSIS
%% ─────────────────────────────────────────────────────────────

\subsection{Prior Sensitivity Analysis}

To assess robustness to prior misspecification, we evaluate RSI under
five prior configurations ranging from domain-informed to deliberately
adversarial:

\begin{center}
\renewcommand{\arraystretch}{1.15}
\begin{tabular}{llcc}
\toprule
\textbf{Configuration} & \textbf{Prior} & $\mathbb{E}[c_i]$ & \textbf{Rationale} \\
\midrule
Correct       & $\mathrm{Beta}(7, 3)$ & 0.70 & Domain-informed (OTR) \\
Uniform       & $\mathrm{Beta}(1, 1)$ & 0.50 & Maximum ignorance \\
Inverted      & $\mathrm{Beta}(3, 7)$ & 0.30 & Deliberately wrong direction \\
Strong wrong  & $\mathrm{Beta}(20, 2)$ & 0.91 & Strong confidence, wrong value \\
Weak          & $\mathrm{Beta}(2, 2)$ & 0.50 & Weak, symmetric \\
\bottomrule
\end{tabular}
\end{center}

Entity-level F1 is perfectly invariant (0.741 across all
configurations) because entity scoring is deterministic:
$\mathrm{NC}_{ji} = 1 - s_{ji}^{\mathrm{raw}}$ does not depend on
the prior. Population posteriors $\mathbb{E}[c_i]$ vary by at most
0.027 across configurations, confirming that with $n > 400$
observations per rule, the likelihood dominates the prior. This
validates the prior robustness corollary of Theorem~\ref{thm:bvm}:
different initial beliefs converge to the same assessment as evidence
accumulates.
\section{Appendix: Supplementary Material}
\label{sec:appendix}

%% ─────────────────────────────────────────────────────────────
%% A.1 — HYPERPARAMETERS
%% ─────────────────────────────────────────────────────────────

\subsection{Hyperparameters}

Table~\ref{tab:hyperparams} lists all hyperparameters used in the RSI
Togolese fiscal instantiation. All values are derived from institutional
knowledge (OTR fiscal rules, historical compliance estimates) or set to
standard non-informative values. No hyperparameter is learned from
labeled data.

\begin{table}[H]
\centering
\caption{RSI hyperparameters. Prior parameters encode institutional
         knowledge; inference parameters are standard.}
\label{tab:hyperparams}
\renewcommand{\arraystretch}{1.15}
\begin{tabular}{llll}
\toprule
\textbf{Parameter} & \textbf{Value} & \textbf{Scope} & \textbf{Justification} \\
\midrule
\multicolumn{4}{l}{\textit{Prior hyperparameters (per rule)}} \\
$\alpha_i, \beta_i$ & See Table~\ref{tab:rule_priors} & Per rule & OTR historical rates \\
$\sigma_i$ (drift) & 0.5--2.0 & Per rule & Expected parametric stability \\
\midrule
\multicolumn{4}{l}{\textit{Inference parameters}} \\
$\sigma_{\mathrm{obs}}$ & 0.0625 & Global & Observation noise scale \\
$\varepsilon$ (ELBO tol.) & $10^{-8}$ & Global & Standard convergence criterion \\
Max iterations & 100 & Global & Upper bound (converges in 2) \\
$\tau$ (discretization) & 0.5 & Global & Neutral threshold \\
\bottomrule
\end{tabular}
\end{table}

\begin{table}[H]
\centering
\caption{Per-rule prior parameters.}
\label{tab:rule_priors}
\renewcommand{\arraystretch}{1.15}
\begin{tabular}{llccl}
\toprule
\textbf{Rule} & \textbf{Description} & $\alpha_i$ & $\beta_i$
  & $\mathbb{E}[c_i]$ \\
\midrule
R1\_TVA  & Value Added Tax          & 7 & 3 & 0.70 \\
R2\_IS   & Corporate Income Tax     & 6 & 4 & 0.60 \\
R3\_IMF  & Minimum Tax              & 7 & 3 & 0.70 \\
R4\_TPU  & Informal Sector Tax      & 5 & 5 & 0.50 \\
R5\_IRPP & Income Tax Withholding   & 6 & 4 & 0.60 \\
R6\_PAT  & Business License         & 7 & 3 & 0.70 \\
R7\_DECL & Annual Declaration       & 6 & 4 & 0.60 \\
R8\_BANK & Bank Formalization       & 5 & 5 & 0.50 \\
\bottomrule
\end{tabular}
\end{table}

%% ─────────────────────────────────────────────────────────────
%% A.2 — DATASET DOCUMENTATION
%% ─────────────────────────────────────────────────────────────

\subsection{RSI-Togo-Fiscal-Synthetic v2.0}

\paragraph{Overview.}
The benchmark contains 2{,}000 synthetic enterprises across two
regulatory periods (2022--2024 and 2025), encoding 8 fiscal rules.
Each enterprise record is structured in four layers: observable
features (segment, declared turnover), latent ground truth
($c_i^*$ per rule), noisy compliance signals with 18\% missing data,
and binary labels for baseline evaluation only.

\paragraph{Market segments.}
Enterprises are assigned to four segments reflecting the structure
of the Togolese market. Turnover follows a log-normal distribution
within each segment, whose aggregate approximates a power
law~\citep{gabaix2009power}.

\begin{center}
\renewcommand{\arraystretch}{1.15}
\begin{tabular}{lcccc}
\toprule
\textbf{Segment} & \textbf{Share} & \textbf{$\mu_{\log}$}
  & \textbf{$\sigma_{\log}$} & \textbf{Median CA (FCFA)} \\
\midrule
Informal      & 45\% & 16.3 & 0.7 & $\sim$12M \\
Small formal  & 25\% & 17.6 & 0.5 & $\sim$44M \\
Medium        & 18\% & 18.6 & 0.4 & $\sim$119M \\
Large         & 12\% & 20.0 & 0.5 & $\sim$485M \\
\bottomrule
\end{tabular}
\end{center}

\paragraph{Under-declaration model.}
Observed turnover is systematically lower than true turnover:
\[
  \hat{x} = x \cdot \underbrace{\beta}_{\text{systematic}} \cdot
  \underbrace{\varepsilon}_{\text{noise}}, \quad
  \beta \sim \operatorname{Beta}(7, 3), \quad
  \varepsilon \sim \mathcal{N}(1,\, 0.04^2)
\]
The Beta(7,3) distribution has mean $0.70$ and concentrates mass
between 0.50 and 0.90, consistent with empirical estimates from
Sub-Saharan African fiscal studies.

\paragraph{Compliance signals.}
For each entity $j$ and applicable rule $r_i$, the continuous
compliance signal is drawn from
$\operatorname{Beta}(c_{ji} \cdot \kappa,\, (1{-}c_{ji}) \cdot \kappa)$
with $\kappa = 10$, ensuring an unbiased signal centered on the true
entity compliance rate: $\mathbb{E}[s_{ji}^{\mathrm{raw}}] = c_{ji}$.
For R4\_TPU (informal sector), the signal is binary:
$s_{ji} \sim \operatorname{Bernoulli}(c_{ji})$.

\paragraph{Applicability and rules.}
Rules R7\_DECL and R8\_BANK apply universally. The remaining 6 rules
have turnover-based thresholds, ensuring every entity is subject to
at least 3 rules (K\,$\geq$\,3). The distribution is:
K\,=\,3 (56\%), K\,=\,5 (15\%), K\,=\,6 (10\%), K\,=\,7 (19\%).

\begin{center}
\renewcommand{\arraystretch}{1.15}
\begin{tabular}{llll}
\toprule
\textbf{Rule} & \textbf{Description} & \textbf{Condition} & \textbf{Threshold} \\
\midrule
R1\_TVA  & Value Added Tax       & CA $\geq$ threshold & 60M (p1), 100M (p2) \\
R2\_IS   & Corporate Income Tax  & CA $\geq$ threshold & 100M \\
R3\_IMF  & Minimum Tax           & CA $\geq$ threshold & 60M \\
R4\_TPU  & Informal Sector Tax   & CA $<$ threshold    & 60M \\
R5\_IRPP & Income Tax Withholding & CA $\geq$ threshold & 30M \\
R6\_PAT  & Business License      & CA $\geq$ threshold & 30M \\
R7\_DECL & Annual Declaration    & Always              & (universal) \\
R8\_BANK & Bank Formalization    & Always              & (universal) \\
\bottomrule
\end{tabular}
\end{center}

\paragraph{Regulatory change event.}
Period~1 (2022--2024) uses a VAT threshold of 60M~FCFA. Period~2
(2025) uses 100M~FCFA, reflecting Law n\textdegree{}2024-007
(30~December~2024). This provides a clean natural experiment for
validating Theorem~\ref{thm:o1}.

\paragraph{Column structure.}
Each enterprise record contains 48 columns: 6 identifiers and
features (\texttt{entity\_id}, \texttt{period}, \texttt{segment},
\texttt{true\_ca}, \texttt{obs\_ca}, \texttt{under\_decl\_ratio}),
5 columns per rule (applicability, true compliance, raw signal,
discretized signal, label), and 2 aggregate columns
(\texttt{label\_any\_nc}, \texttt{n\_applicable}).

%% ─────────────────────────────────────────────────────────────
%% A.3 — BASELINES
%% ─────────────────────────────────────────────────────────────

\subsection{Baseline Configurations}

\begin{table}[H]
\centering
\caption{Baseline model configurations.}
\label{tab:baselines}
\renewcommand{\arraystretch}{1.15}
\begin{tabular}{lll}
\toprule
\textbf{Model} & \textbf{Parameter} & \textbf{Value} \\
\midrule
LR      & solver & lbfgs \\
        & max\_iter & 1000 \\
        & feature scaling & StandardScaler \\
        & train/test split & 70\%/30\%, stratified \\[4pt]
XGBoost & n\_estimators & 150 \\
        & max\_depth & 4 \\
        & learning\_rate & 0.1 (default) \\
        & train/test split & 70\%/30\%, stratified \\[4pt]
MLP     & hidden\_layer\_sizes & (128, 64, 32) \\
        & activation & ReLU \\
        & max\_iter & 1000 \\
        & early\_stopping & True \\
        & learning\_rate & adaptive \\
        & feature scaling & StandardScaler \\[4pt]
RBS     & Logic & Flag if $s_{ji}^{\mathrm{disc}} = 0$ for any applicable rule \\
        & Missing handling & Ignore (no flag) \\
\bottomrule
\end{tabular}
\end{table}

%% ─────────────────────────────────────────────────────────────
%% A.4 — REPRODUCIBILITY
%% ─────────────────────────────────────────────────────────────

\subsection{Reproducibility}

All experiments use a fixed random seed (42) for full deterministic
reproducibility. The implementation uses Python~3.11 with NumPy,
SciPy, scikit-learn, and pandas. No GPU is required. The full
experiment pipeline completes in under 30~seconds on standard
hardware. The dataset, framework, and all experiment code are
released publicly under the MIT license.

%% ─────────────────────────────────────────────────────────────
%% A.5 — DETAILED PROOFS
%% ─────────────────────────────────────────────────────────────

\subsection{Detailed Proof of Theorem~\ref{thm:o1} ($O(1)$ Adaptability)}

\begin{proof}
Let $\mathcal{U}_k$ be a regulatory update that changes the parameters
of rule $r_k$ from $\Theta_k$ to $\Theta_k'$.

\textit{Step 1: Prior factorization.}
The RSI prior factorizes as
$P(\mathbf{S}) = \prod_{i=1}^n P(c_i \mid \Theta_i) \cdot P(\delta_i \mid \Theta_i)$.
The posterior satisfies:
\[
  P(\mathbf{S} \mid \mathbf{D}) \propto
  \prod_{i=1}^n P(\mathbf{D}_i \mid c_i) \cdot P(c_i \mid \Theta_i) \cdot P(\delta_i \mid \Theta_i)
\]
After the update $\Theta_k \to \Theta_k'$, only the $k$-th factor
changes:
\[
  P'(\mathbf{S} \mid \mathbf{D}) \propto
  P(\mathbf{S} \mid \mathbf{D}) \cdot
  \frac{P(c_k \mid \Theta_k') \cdot P(\delta_k \mid \Theta_k')}
       {P(c_k \mid \Theta_k) \cdot P(\delta_k \mid \Theta_k)}
\]

\textit{Step 2: Drift correction.}
For a threshold change $\theta_k \to \theta_k'$, define
$\varepsilon = (\theta_k' - \theta_k)/\theta_k$. The drift posterior
mean updates as $\bar{\mu}_k \leftarrow \bar{\mu}_k - \varepsilon$.
This is a single scalar subtraction: $O(1)$.

\textit{Step 3: Invariance of compliance parameters.}
The compliance rate $c_k$ is independent of the threshold magnitude
(it measures the \emph{proportion} of applicable entities that comply,
not the threshold itself). Therefore $\bar{\alpha}_k$ and
$\bar{\beta}_k$ are invariant under $\mathcal{U}_k$. Similarly,
$\bar{\sigma}_k$ depends on the number of observations and prior
variance, neither of which changes under $\mathcal{U}_k$.

\textit{Step 4: Normalization.}
Under conjugate families, the posterior normalization constant
$B(\bar{\alpha}_k, \bar{\beta}_k)$ is computed analytically and does
not require re-processing the data. The updated posterior is a
properly normalized Beta distribution.

\textit{Total cost:} one scalar subtraction ($\bar{\mu}_k$), zero
data access. $\mathcal{C}^{\mathrm{RSI}}(\mathcal{U}_k) = O(1)$
with respect to $|\mathbf{D}|$ and $n$. \qed
\end{proof}

\subsection{Detailed Proof of Theorem~\ref{thm:bvm} (BvM Consistency)}

\begin{proof}
We prove the result for a single rule $r_i$ (the extension to
$\mathbf{S}$ follows by factorization).

Let $s_1, s_2, \ldots, s_m$ be i.i.d.\ Bernoulli($c_i^*$) compliance
signals, where $c_i^* \in (0,1)$ is the true population compliance rate.

\textit{Step 1: Posterior computation.}
The conjugate posterior is
$Q^*(c_i) = \operatorname{Beta}(\bar{\alpha}_i, \bar{\beta}_i)$ with
$\bar{\alpha}_i = \alpha_i + \sum_{j=1}^m s_j$ and
$\bar{\beta}_i = \beta_i + \sum_{j=1}^m (1-s_j)$.

\textit{Step 2: Posterior mean convergence.}
Define $\bar{s} = m^{-1}\sum_j s_j$. Then:
\[
  \mathbb{E}_Q[c_i] = \frac{\bar{\alpha}_i}{\bar{\alpha}_i + \bar{\beta}_i}
  = \frac{\alpha_i + m\bar{s}}{\alpha_i + \beta_i + m}
  \to \bar{s} \to c_i^* \quad (m \to \infty)
\]
where the first convergence follows from $m \gg \alpha_i + \beta_i$
and the second from the strong law of large numbers.

\textit{Step 3: Posterior variance.}
\[
  \mathrm{Var}_Q[c_i] = \frac{\bar{\alpha}_i \bar{\beta}_i}
  {(\bar{\alpha}_i + \bar{\beta}_i)^2 (\bar{\alpha}_i + \bar{\beta}_i + 1)}
  \sim \frac{c_i^*(1-c_i^*)}{m} \quad (m \to \infty)
\]
This is the Bernoulli variance divided by $m$, confirming
$\sigma[c_i] \sim 1/\sqrt{m}$.

\textit{Step 4: Asymptotic normality.}
By the central limit theorem, $\sqrt{m}(\bar{s} - c_i^*)
\xrightarrow{d} \mathcal{N}(0, c_i^*(1-c_i^*))$. Since
$\mathbb{E}_Q[c_i] \to \bar{s}$ and $\mathrm{Var}_Q[c_i] \to
c_i^*(1-c_i^*)/m$, we obtain:
\[
  \sqrt{m}(\mathbb{E}_Q[c_i] - c_i^*) \xrightarrow{d}
  \mathcal{N}(0, \mathcal{I}(c_i^*)^{-1})
\]
where $\mathcal{I}(c_i^*) = [c_i^*(1-c_i^*)]^{-1}$ is the Fisher
information of the Bernoulli model. This is the classical
Bernstein-von Mises theorem applied to the Beta-Bernoulli conjugate
pair \citep{van2000asymptotic}. \qed
\end{proof}

%% ─────────────────────────────────────────────────────────────
%% A.6 — SIGNAL DISCRETIZATION BIAS
%% ─────────────────────────────────────────────────────────────

\subsection{Signal Discretization and Calibration Bias}

The compliance signal $s_{ji}^{\mathrm{raw}} \in [0,1]$ is unbiased
by construction:
$\mathbb{E}[s_{ji}^{\mathrm{raw}}] = c_{ji}$.
Discretization at threshold $\tau = 0.5$ introduces a systematic bias.
For an entity with true compliance $c_{ji}$, the probability that the
discretized signal equals 1 is:
\[
  P(s_{ji}^{\mathrm{disc}} = 1) =
  1 - I_{0.5}(c_{ji} \cdot \kappa,\, (1{-}c_{ji}) \cdot \kappa)
\]
where $I_x(a,b)$ is the regularized incomplete Beta function. When
$c_{ji} > 0.5$, this probability exceeds $c_{ji}$ (the discrete signal
overestimates compliance). When $c_{ji} < 0.5$, it underestimates.
The bias is largest for $c_{ji}$ far from 0.5 and decreases as
$\kappa \to \infty$.

This explains the calibration misses in Table~\ref{tab:calibration}:
rules with true population compliance far from 0.5 (R3\_IMF at 0.876,
R7\_DECL at 0.666) exhibit the largest bias, while rules near 0.5
(R4\_TPU at 0.345, R8\_BANK at 0.491) are well-calibrated.

A natural mitigation is to use continuous signals directly with a
Beta-likelihood formulation, eliminating discretization entirely while
preserving conjugacy under moment-matching approximations. This
extension is developed in future work.

%% ─────────────────────────────────────────────────────────────
%% A.7 — PRIOR SENSITIVITY ANALYSIS
%% ─────────────────────────────────────────────────────────────

\subsection{Prior Sensitivity Analysis}

To assess robustness to prior misspecification, we evaluate RSI under
five prior configurations ranging from domain-informed to deliberately
adversarial:

\begin{center}
\renewcommand{\arraystretch}{1.15}
\begin{tabular}{llcc}
\toprule
\textbf{Configuration} & \textbf{Prior} & $\mathbb{E}[c_i]$ & \textbf{Rationale} \\
\midrule
Correct       & $\mathrm{Beta}(7, 3)$ & 0.70 & Domain-informed (OTR) \\
Uniform       & $\mathrm{Beta}(1, 1)$ & 0.50 & Maximum ignorance \\
Inverted      & $\mathrm{Beta}(3, 7)$ & 0.30 & Deliberately wrong direction \\
Strong wrong  & $\mathrm{Beta}(20, 2)$ & 0.91 & Strong confidence, wrong value \\
Weak          & $\mathrm{Beta}(2, 2)$ & 0.50 & Weak, symmetric \\
\bottomrule
\end{tabular}
\end{center}

Entity-level F1 is perfectly invariant (0.741 across all
configurations) because entity scoring is deterministic:
$\mathrm{NC}_{ji} = 1 - s_{ji}^{\mathrm{raw}}$ does not depend on
the prior. Population posteriors $\mathbb{E}[c_i]$ vary by at most
0.027 across configurations, confirming that with $n > 400$
observations per rule, the likelihood dominates the prior. This
validates the prior robustness corollary of Theorem~\ref{thm:bvm}:
different initial beliefs converge to the same assessment as evidence
accumulates.
\section{Appendix: Supplementary Material}
\label{sec:appendix}

%% ─────────────────────────────────────────────────────────────
%% A.1 — HYPERPARAMETERS
%% ─────────────────────────────────────────────────────────────

\subsection{Hyperparameters}

Table~\ref{tab:hyperparams} lists all hyperparameters used in the RSI
Togolese fiscal instantiation. All values are derived from institutional
knowledge (OTR fiscal rules, historical compliance estimates) or set to
standard non-informative values. No hyperparameter is learned from
labeled data.

\begin{table}[H]
\centering
\caption{RSI hyperparameters. Prior parameters encode institutional
         knowledge; inference parameters are standard.}
\label{tab:hyperparams}
\renewcommand{\arraystretch}{1.15}
\begin{tabular}{llll}
\toprule
\textbf{Parameter} & \textbf{Value} & \textbf{Scope} & \textbf{Justification} \\
\midrule
\multicolumn{4}{l}{\textit{Prior hyperparameters (per rule)}} \\
$\alpha_i, \beta_i$ & See Table~\ref{tab:rule_priors} & Per rule & OTR historical rates \\
$\sigma_i$ (drift) & 0.5--2.0 & Per rule & Expected parametric stability \\
\midrule
\multicolumn{4}{l}{\textit{Inference parameters}} \\
$\sigma_{\mathrm{obs}}$ & 0.0625 & Global & Observation noise scale \\
$\varepsilon$ (ELBO tol.) & $10^{-8}$ & Global & Standard convergence criterion \\
Max iterations & 100 & Global & Upper bound (converges in 2) \\
$\tau$ (discretization) & 0.5 & Global & Neutral threshold \\
\bottomrule
\end{tabular}
\end{table}

\begin{table}[H]
\centering
\caption{Per-rule prior parameters.}
\label{tab:rule_priors}
\renewcommand{\arraystretch}{1.15}
\begin{tabular}{llccl}
\toprule
\textbf{Rule} & \textbf{Description} & $\alpha_i$ & $\beta_i$
  & $\mathbb{E}[c_i]$ \\
\midrule
R1\_TVA  & Value Added Tax          & 7 & 3 & 0.70 \\
R2\_IS   & Corporate Income Tax     & 6 & 4 & 0.60 \\
R3\_IMF  & Minimum Tax              & 7 & 3 & 0.70 \\
R4\_TPU  & Informal Sector Tax      & 5 & 5 & 0.50 \\
R5\_IRPP & Income Tax Withholding   & 6 & 4 & 0.60 \\
R6\_PAT  & Business License         & 7 & 3 & 0.70 \\
R7\_DECL & Annual Declaration       & 6 & 4 & 0.60 \\
R8\_BANK & Bank Formalization       & 5 & 5 & 0.50 \\
\bottomrule
\end{tabular}
\end{table}

%% ─────────────────────────────────────────────────────────────
%% A.2 — DATASET DOCUMENTATION
%% ─────────────────────────────────────────────────────────────

\subsection{RSI-Togo-Fiscal-Synthetic v2.0}

\paragraph{Overview.}
The benchmark contains 2{,}000 synthetic enterprises across two
regulatory periods (2022--2024 and 2025), encoding 8 fiscal rules.
Each enterprise record is structured in four layers: observable
features (segment, declared turnover), latent ground truth
($c_i^*$ per rule), noisy compliance signals with 18\% missing data,
and binary labels for baseline evaluation only.

\paragraph{Market segments.}
Enterprises are assigned to four segments reflecting the structure
of the Togolese market. Turnover follows a log-normal distribution
within each segment, whose aggregate approximates a power
law~\citep{gabaix2009power}.

\begin{center}
\renewcommand{\arraystretch}{1.15}
\begin{tabular}{lcccc}
\toprule
\textbf{Segment} & \textbf{Share} & \textbf{$\mu_{\log}$}
  & \textbf{$\sigma_{\log}$} & \textbf{Median CA (FCFA)} \\
\midrule
Informal      & 45\% & 16.3 & 0.7 & $\sim$12M \\
Small formal  & 25\% & 17.6 & 0.5 & $\sim$44M \\
Medium        & 18\% & 18.6 & 0.4 & $\sim$119M \\
Large         & 12\% & 20.0 & 0.5 & $\sim$485M \\
\bottomrule
\end{tabular}
\end{center}

\paragraph{Under-declaration model.}
Observed turnover is systematically lower than true turnover:
\[
  \hat{x} = x \cdot \underbrace{\beta}_{\text{systematic}} \cdot
  \underbrace{\varepsilon}_{\text{noise}}, \quad
  \beta \sim \operatorname{Beta}(7, 3), \quad
  \varepsilon \sim \mathcal{N}(1,\, 0.04^2)
\]
The Beta(7,3) distribution has mean $0.70$ and concentrates mass
between 0.50 and 0.90, consistent with empirical estimates from
Sub-Saharan African fiscal studies.

\paragraph{Compliance signals.}
For each entity $j$ and applicable rule $r_i$, the continuous
compliance signal is drawn from
$\operatorname{Beta}(c_{ji} \cdot \kappa,\, (1{-}c_{ji}) \cdot \kappa)$
with $\kappa = 10$, ensuring an unbiased signal centered on the true
entity compliance rate: $\mathbb{E}[s_{ji}^{\mathrm{raw}}] = c_{ji}$.
For R4\_TPU (informal sector), the signal is binary:
$s_{ji} \sim \operatorname{Bernoulli}(c_{ji})$.

\paragraph{Applicability and rules.}
Rules R7\_DECL and R8\_BANK apply universally. The remaining 6 rules
have turnover-based thresholds, ensuring every entity is subject to
at least 3 rules (K\,$\geq$\,3). The distribution is:
K\,=\,3 (56\%), K\,=\,5 (15\%), K\,=\,6 (10\%), K\,=\,7 (19\%).

\begin{center}
\renewcommand{\arraystretch}{1.15}
\begin{tabular}{llll}
\toprule
\textbf{Rule} & \textbf{Description} & \textbf{Condition} & \textbf{Threshold} \\
\midrule
R1\_TVA  & Value Added Tax       & CA $\geq$ threshold & 60M (p1), 100M (p2) \\
R2\_IS   & Corporate Income Tax  & CA $\geq$ threshold & 100M \\
R3\_IMF  & Minimum Tax           & CA $\geq$ threshold & 60M \\
R4\_TPU  & Informal Sector Tax   & CA $<$ threshold    & 60M \\
R5\_IRPP & Income Tax Withholding & CA $\geq$ threshold & 30M \\
R6\_PAT  & Business License      & CA $\geq$ threshold & 30M \\
R7\_DECL & Annual Declaration    & Always              & (universal) \\
R8\_BANK & Bank Formalization    & Always              & (universal) \\
\bottomrule
\end{tabular}
\end{center}

\paragraph{Regulatory change event.}
Period~1 (2022--2024) uses a VAT threshold of 60M~FCFA. Period~2
(2025) uses 100M~FCFA, reflecting Law n\textdegree{}2024-007
(30~December~2024). This provides a clean natural experiment for
validating Theorem~\ref{thm:o1}.

\paragraph{Column structure.}
Each enterprise record contains 48 columns: 6 identifiers and
features (\texttt{entity\_id}, \texttt{period}, \texttt{segment},
\texttt{true\_ca}, \texttt{obs\_ca}, \texttt{under\_decl\_ratio}),
5 columns per rule (applicability, true compliance, raw signal,
discretized signal, label), and 2 aggregate columns
(\texttt{label\_any\_nc}, \texttt{n\_applicable}).

%% ─────────────────────────────────────────────────────────────
%% A.3 — BASELINES
%% ─────────────────────────────────────────────────────────────

\subsection{Baseline Configurations}

\begin{table}[H]
\centering
\caption{Baseline model configurations.}
\label{tab:baselines}
\renewcommand{\arraystretch}{1.15}
\begin{tabular}{lll}
\toprule
\textbf{Model} & \textbf{Parameter} & \textbf{Value} \\
\midrule
LR      & solver & lbfgs \\
        & max\_iter & 1000 \\
        & feature scaling & StandardScaler \\
        & train/test split & 70\%/30\%, stratified \\[4pt]
XGBoost & n\_estimators & 150 \\
        & max\_depth & 4 \\
        & learning\_rate & 0.1 (default) \\
        & train/test split & 70\%/30\%, stratified \\[4pt]
MLP     & hidden\_layer\_sizes & (128, 64, 32) \\
        & activation & ReLU \\
        & max\_iter & 1000 \\
        & early\_stopping & True \\
        & learning\_rate & adaptive \\
        & feature scaling & StandardScaler \\[4pt]
RBS     & Logic & Flag if $s_{ji}^{\mathrm{disc}} = 0$ for any applicable rule \\
        & Missing handling & Ignore (no flag) \\
\bottomrule
\end{tabular}
\end{table}

%% ─────────────────────────────────────────────────────────────
%% A.4 — REPRODUCIBILITY
%% ─────────────────────────────────────────────────────────────

\subsection{Reproducibility}

All experiments use a fixed random seed (42) for full deterministic
reproducibility. The implementation uses Python~3.11 with NumPy,
SciPy, scikit-learn, and pandas. No GPU is required. The full
experiment pipeline completes in under 30~seconds on standard
hardware. The dataset, framework, and all experiment code are
released publicly under the MIT license.

%% ─────────────────────────────────────────────────────────────
%% A.5 — DETAILED PROOFS
%% ─────────────────────────────────────────────────────────────

\subsection{Detailed Proof of Theorem~\ref{thm:o1} ($O(1)$ Adaptability)}

\begin{proof}
Let $\mathcal{U}_k$ be a regulatory update that changes the parameters
of rule $r_k$ from $\Theta_k$ to $\Theta_k'$.

\textit{Step 1: Prior factorization.}
The RSI prior factorizes as
$P(\mathbf{S}) = \prod_{i=1}^n P(c_i \mid \Theta_i) \cdot P(\delta_i \mid \Theta_i)$.
The posterior satisfies:
\[
  P(\mathbf{S} \mid \mathbf{D}) \propto
  \prod_{i=1}^n P(\mathbf{D}_i \mid c_i) \cdot P(c_i \mid \Theta_i) \cdot P(\delta_i \mid \Theta_i)
\]
After the update $\Theta_k \to \Theta_k'$, only the $k$-th factor
changes:
\[
  P'(\mathbf{S} \mid \mathbf{D}) \propto
  P(\mathbf{S} \mid \mathbf{D}) \cdot
  \frac{P(c_k \mid \Theta_k') \cdot P(\delta_k \mid \Theta_k')}
       {P(c_k \mid \Theta_k) \cdot P(\delta_k \mid \Theta_k)}
\]

\textit{Step 2: Drift correction.}
For a threshold change $\theta_k \to \theta_k'$, define
$\varepsilon = (\theta_k' - \theta_k)/\theta_k$. The drift posterior
mean updates as $\bar{\mu}_k \leftarrow \bar{\mu}_k - \varepsilon$.
This is a single scalar subtraction: $O(1)$.

\textit{Step 3: Invariance of compliance parameters.}
The compliance rate $c_k$ is independent of the threshold magnitude
(it measures the \emph{proportion} of applicable entities that comply,
not the threshold itself). Therefore $\bar{\alpha}_k$ and
$\bar{\beta}_k$ are invariant under $\mathcal{U}_k$. Similarly,
$\bar{\sigma}_k$ depends on the number of observations and prior
variance, neither of which changes under $\mathcal{U}_k$.

\textit{Step 4: Normalization.}
Under conjugate families, the posterior normalization constant
$B(\bar{\alpha}_k, \bar{\beta}_k)$ is computed analytically and does
not require re-processing the data. The updated posterior is a
properly normalized Beta distribution.

\textit{Total cost:} one scalar subtraction ($\bar{\mu}_k$), zero
data access. $\mathcal{C}^{\mathrm{RSI}}(\mathcal{U}_k) = O(1)$
with respect to $|\mathbf{D}|$ and $n$. \qed
\end{proof}

\subsection{Detailed Proof of Theorem~\ref{thm:bvm} (BvM Consistency)}

\begin{proof}
We prove the result for a single rule $r_i$ (the extension to
$\mathbf{S}$ follows by factorization).

Let $s_1, s_2, \ldots, s_m$ be i.i.d.\ Bernoulli($c_i^*$) compliance
signals, where $c_i^* \in (0,1)$ is the true population compliance rate.

\textit{Step 1: Posterior computation.}
The conjugate posterior is
$Q^*(c_i) = \operatorname{Beta}(\bar{\alpha}_i, \bar{\beta}_i)$ with
$\bar{\alpha}_i = \alpha_i + \sum_{j=1}^m s_j$ and
$\bar{\beta}_i = \beta_i + \sum_{j=1}^m (1-s_j)$.

\textit{Step 2: Posterior mean convergence.}
Define $\bar{s} = m^{-1}\sum_j s_j$. Then:
\[
  \mathbb{E}_Q[c_i] = \frac{\bar{\alpha}_i}{\bar{\alpha}_i + \bar{\beta}_i}
  = \frac{\alpha_i + m\bar{s}}{\alpha_i + \beta_i + m}
  \to \bar{s} \to c_i^* \quad (m \to \infty)
\]
where the first convergence follows from $m \gg \alpha_i + \beta_i$
and the second from the strong law of large numbers.

\textit{Step 3: Posterior variance.}
\[
  \mathrm{Var}_Q[c_i] = \frac{\bar{\alpha}_i \bar{\beta}_i}
  {(\bar{\alpha}_i + \bar{\beta}_i)^2 (\bar{\alpha}_i + \bar{\beta}_i + 1)}
  \sim \frac{c_i^*(1-c_i^*)}{m} \quad (m \to \infty)
\]
This is the Bernoulli variance divided by $m$, confirming
$\sigma[c_i] \sim 1/\sqrt{m}$.

\textit{Step 4: Asymptotic normality.}
By the central limit theorem, $\sqrt{m}(\bar{s} - c_i^*)
\xrightarrow{d} \mathcal{N}(0, c_i^*(1-c_i^*))$. Since
$\mathbb{E}_Q[c_i] \to \bar{s}$ and $\mathrm{Var}_Q[c_i] \to
c_i^*(1-c_i^*)/m$, we obtain:
\[
  \sqrt{m}(\mathbb{E}_Q[c_i] - c_i^*) \xrightarrow{d}
  \mathcal{N}(0, \mathcal{I}(c_i^*)^{-1})
\]
where $\mathcal{I}(c_i^*) = [c_i^*(1-c_i^*)]^{-1}$ is the Fisher
information of the Bernoulli model. This is the classical
Bernstein-von Mises theorem applied to the Beta-Bernoulli conjugate
pair \citep{van2000asymptotic}. \qed
\end{proof}

%% ─────────────────────────────────────────────────────────────
%% A.6 — SIGNAL DISCRETIZATION BIAS
%% ─────────────────────────────────────────────────────────────

\subsection{Signal Discretization and Calibration Bias}

The compliance signal $s_{ji}^{\mathrm{raw}} \in [0,1]$ is unbiased
by construction:
$\mathbb{E}[s_{ji}^{\mathrm{raw}}] = c_{ji}$.
Discretization at threshold $\tau = 0.5$ introduces a systematic bias.
For an entity with true compliance $c_{ji}$, the probability that the
discretized signal equals 1 is:
\[
  P(s_{ji}^{\mathrm{disc}} = 1) =
  1 - I_{0.5}(c_{ji} \cdot \kappa,\, (1{-}c_{ji}) \cdot \kappa)
\]
where $I_x(a,b)$ is the regularized incomplete Beta function. When
$c_{ji} > 0.5$, this probability exceeds $c_{ji}$ (the discrete signal
overestimates compliance). When $c_{ji} < 0.5$, it underestimates.
The bias is largest for $c_{ji}$ far from 0.5 and decreases as
$\kappa \to \infty$.

This explains the calibration misses in Table~\ref{tab:calibration}:
rules with true population compliance far from 0.5 (R3\_IMF at 0.876,
R7\_DECL at 0.666) exhibit the largest bias, while rules near 0.5
(R4\_TPU at 0.345, R8\_BANK at 0.491) are well-calibrated.

A natural mitigation is to use continuous signals directly with a
Beta-likelihood formulation, eliminating discretization entirely while
preserving conjugacy under moment-matching approximations. This
extension is developed in future work.

%% ─────────────────────────────────────────────────────────────
%% A.7 — PRIOR SENSITIVITY ANALYSIS
%% ─────────────────────────────────────────────────────────────

\subsection{Prior Sensitivity Analysis}

To assess robustness to prior misspecification, we evaluate RSI under
five prior configurations ranging from domain-informed to deliberately
adversarial:

\begin{center}
\renewcommand{\arraystretch}{1.15}
\begin{tabular}{llcc}
\toprule
\textbf{Configuration} & \textbf{Prior} & $\mathbb{E}[c_i]$ & \textbf{Rationale} \\
\midrule
Correct       & $\mathrm{Beta}(7, 3)$ & 0.70 & Domain-informed (OTR) \\
Uniform       & $\mathrm{Beta}(1, 1)$ & 0.50 & Maximum ignorance \\
Inverted      & $\mathrm{Beta}(3, 7)$ & 0.30 & Deliberately wrong direction \\
Strong wrong  & $\mathrm{Beta}(20, 2)$ & 0.91 & Strong confidence, wrong value \\
Weak          & $\mathrm{Beta}(2, 2)$ & 0.50 & Weak, symmetric \\
\bottomrule
\end{tabular}
\end{center}

Entity-level F1 is perfectly invariant (0.741 across all
configurations) because entity scoring is deterministic:
$\mathrm{NC}_{ji} = 1 - s_{ji}^{\mathrm{raw}}$ does not depend on
the prior. Population posteriors $\mathbb{E}[c_i]$ vary by at most
0.027 across configurations, confirming that with $n > 400$
observations per rule, the likelihood dominates the prior. This
validates the prior robustness corollary of Theorem~\ref{thm:bvm}:
different initial beliefs converge to the same assessment as evidence
accumulates.
\section{Appendix: Supplementary Material}
\label{sec:appendix}

%% ─────────────────────────────────────────────────────────────
%% A.1 — HYPERPARAMETERS
%% ─────────────────────────────────────────────────────────────

\subsection{Hyperparameters}

Table~\ref{tab:hyperparams} lists all hyperparameters used in the RSI
Togolese fiscal instantiation. All values are derived from institutional
knowledge (OTR fiscal rules, historical compliance estimates) or set to
standard non-informative values. No hyperparameter is learned from
labeled data.

\begin{table}[H]
\centering
\caption{RSI hyperparameters. Prior parameters encode institutional
         knowledge; inference parameters are standard.}
\label{tab:hyperparams}
\renewcommand{\arraystretch}{1.15}
\begin{tabular}{llll}
\toprule
\textbf{Parameter} & \textbf{Value} & \textbf{Scope} & \textbf{Justification} \\
\midrule
\multicolumn{4}{l}{\textit{Prior hyperparameters (per rule)}} \\
$\alpha_i, \beta_i$ & See Table~\ref{tab:rule_priors} & Per rule & OTR historical rates \\
$\sigma_i$ (drift) & 0.5--2.0 & Per rule & Expected parametric stability \\
\midrule
\multicolumn{4}{l}{\textit{Inference parameters}} \\
$\sigma_{\mathrm{obs}}$ & 0.0625 & Global & Observation noise scale \\
$\varepsilon$ (ELBO tol.) & $10^{-8}$ & Global & Standard convergence criterion \\
Max iterations & 100 & Global & Upper bound (converges in 2) \\
$\tau$ (discretization) & 0.5 & Global & Neutral threshold \\
\bottomrule
\end{tabular}
\end{table}

\begin{table}[H]
\centering
\caption{Per-rule prior parameters.}
\label{tab:rule_priors}
\renewcommand{\arraystretch}{1.15}
\begin{tabular}{llccl}
\toprule
\textbf{Rule} & \textbf{Description} & $\alpha_i$ & $\beta_i$
  & $\mathbb{E}[c_i]$ \\
\midrule
R1\_TVA  & Value Added Tax          & 7 & 3 & 0.70 \\
R2\_IS   & Corporate Income Tax     & 6 & 4 & 0.60 \\
R3\_IMF  & Minimum Tax              & 7 & 3 & 0.70 \\
R4\_TPU  & Informal Sector Tax      & 5 & 5 & 0.50 \\
R5\_IRPP & Income Tax Withholding   & 6 & 4 & 0.60 \\
R6\_PAT  & Business License         & 7 & 3 & 0.70 \\
R7\_DECL & Annual Declaration       & 6 & 4 & 0.60 \\
R8\_BANK & Bank Formalization       & 5 & 5 & 0.50 \\
\bottomrule
\end{tabular}
\end{table}

%% ─────────────────────────────────────────────────────────────
%% A.2 — DATASET DOCUMENTATION
%% ─────────────────────────────────────────────────────────────

\subsection{RSI-Togo-Fiscal-Synthetic v2.0}

\paragraph{Overview.}
The benchmark contains 2{,}000 synthetic enterprises across two
regulatory periods (2022--2024 and 2025), encoding 8 fiscal rules.
Each enterprise record is structured in four layers: observable
features (segment, declared turnover), latent ground truth
($c_i^*$ per rule), noisy compliance signals with 18\% missing data,
and binary labels for baseline evaluation only.

\paragraph{Market segments.}
Enterprises are assigned to four segments reflecting the structure
of the Togolese market. Turnover follows a log-normal distribution
within each segment, whose aggregate approximates a power
law~\citep{gabaix2009power}.

\begin{center}
\renewcommand{\arraystretch}{1.15}
\begin{tabular}{lcccc}
\toprule
\textbf{Segment} & \textbf{Share} & \textbf{$\mu_{\log}$}
  & \textbf{$\sigma_{\log}$} & \textbf{Median CA (FCFA)} \\
\midrule
Informal      & 45\% & 16.3 & 0.7 & $\sim$12M \\
Small formal  & 25\% & 17.6 & 0.5 & $\sim$44M \\
Medium        & 18\% & 18.6 & 0.4 & $\sim$119M \\
Large         & 12\% & 20.0 & 0.5 & $\sim$485M \\
\bottomrule
\end{tabular}
\end{center}

\paragraph{Under-declaration model.}
Observed turnover is systematically lower than true turnover:
\[
  \hat{x} = x \cdot \underbrace{\beta}_{\text{systematic}} \cdot
  \underbrace{\varepsilon}_{\text{noise}}, \quad
  \beta \sim \operatorname{Beta}(7, 3), \quad
  \varepsilon \sim \mathcal{N}(1,\, 0.04^2)
\]
The Beta(7,3) distribution has mean $0.70$ and concentrates mass
between 0.50 and 0.90, consistent with empirical estimates from
Sub-Saharan African fiscal studies.

\paragraph{Compliance signals.}
For each entity $j$ and applicable rule $r_i$, the continuous
compliance signal is drawn from
$\operatorname{Beta}(c_{ji} \cdot \kappa,\, (1{-}c_{ji}) \cdot \kappa)$
with $\kappa = 10$, ensuring an unbiased signal centered on the true
entity compliance rate: $\mathbb{E}[s_{ji}^{\mathrm{raw}}] = c_{ji}$.
For R4\_TPU (informal sector), the signal is binary:
$s_{ji} \sim \operatorname{Bernoulli}(c_{ji})$.

\paragraph{Applicability and rules.}
Rules R7\_DECL and R8\_BANK apply universally. The remaining 6 rules
have turnover-based thresholds, ensuring every entity is subject to
at least 3 rules (K\,$\geq$\,3). The distribution is:
K\,=\,3 (56\%), K\,=\,5 (15\%), K\,=\,6 (10\%), K\,=\,7 (19\%).

\begin{center}
\renewcommand{\arraystretch}{1.15}
\begin{tabular}{llll}
\toprule
\textbf{Rule} & \textbf{Description} & \textbf{Condition} & \textbf{Threshold} \\
\midrule
R1\_TVA  & Value Added Tax       & CA $\geq$ threshold & 60M (p1), 100M (p2) \\
R2\_IS   & Corporate Income Tax  & CA $\geq$ threshold & 100M \\
R3\_IMF  & Minimum Tax           & CA $\geq$ threshold & 60M \\
R4\_TPU  & Informal Sector Tax   & CA $<$ threshold    & 60M \\
R5\_IRPP & Income Tax Withholding & CA $\geq$ threshold & 30M \\
R6\_PAT  & Business License      & CA $\geq$ threshold & 30M \\
R7\_DECL & Annual Declaration    & Always              & (universal) \\
R8\_BANK & Bank Formalization    & Always              & (universal) \\
\bottomrule
\end{tabular}
\end{center}

\paragraph{Regulatory change event.}
Period~1 (2022--2024) uses a VAT threshold of 60M~FCFA. Period~2
(2025) uses 100M~FCFA, reflecting Law n\textdegree{}2024-007
(30~December~2024). This provides a clean natural experiment for
validating Theorem~\ref{thm:o1}.

\paragraph{Column structure.}
Each enterprise record contains 48 columns: 6 identifiers and
features (\texttt{entity\_id}, \texttt{period}, \texttt{segment},
\texttt{true\_ca}, \texttt{obs\_ca}, \texttt{under\_decl\_ratio}),
5 columns per rule (applicability, true compliance, raw signal,
discretized signal, label), and 2 aggregate columns
(\texttt{label\_any\_nc}, \texttt{n\_applicable}).

%% ─────────────────────────────────────────────────────────────
%% A.3 — BASELINES
%% ─────────────────────────────────────────────────────────────

\subsection{Baseline Configurations}

\begin{table}[H]
\centering
\caption{Baseline model configurations.}
\label{tab:baselines}
\renewcommand{\arraystretch}{1.15}
\begin{tabular}{lll}
\toprule
\textbf{Model} & \textbf{Parameter} & \textbf{Value} \\
\midrule
LR      & solver & lbfgs \\
        & max\_iter & 1000 \\
        & feature scaling & StandardScaler \\
        & train/test split & 70\%/30\%, stratified \\[4pt]
XGBoost & n\_estimators & 150 \\
        & max\_depth & 4 \\
        & learning\_rate & 0.1 (default) \\
        & train/test split & 70\%/30\%, stratified \\[4pt]
MLP     & hidden\_layer\_sizes & (128, 64, 32) \\
        & activation & ReLU \\
        & max\_iter & 1000 \\
        & early\_stopping & True \\
        & learning\_rate & adaptive \\
        & feature scaling & StandardScaler \\[4pt]
RBS     & Logic & Flag if $s_{ji}^{\mathrm{disc}} = 0$ for any applicable rule \\
        & Missing handling & Ignore (no flag) \\
\bottomrule
\end{tabular}
\end{table}

%% ─────────────────────────────────────────────────────────────
%% A.4 — REPRODUCIBILITY
%% ─────────────────────────────────────────────────────────────

\subsection{Reproducibility}

All experiments use a fixed random seed (42) for full deterministic
reproducibility. The implementation uses Python~3.11 with NumPy,
SciPy, scikit-learn, and pandas. No GPU is required. The full
experiment pipeline completes in under 30~seconds on standard
hardware. The dataset, framework, and all experiment code are
released publicly under the MIT license.

%% ─────────────────────────────────────────────────────────────
%% A.5 — DETAILED PROOFS
%% ─────────────────────────────────────────────────────────────

\subsection{Detailed Proof of Theorem~\ref{thm:o1} ($O(1)$ Adaptability)}

\begin{proof}
Let $\mathcal{U}_k$ be a regulatory update that changes the parameters
of rule $r_k$ from $\Theta_k$ to $\Theta_k'$.

\textit{Step 1: Prior factorization.}
The RSI prior factorizes as
$P(\mathbf{S}) = \prod_{i=1}^n P(c_i \mid \Theta_i) \cdot P(\delta_i \mid \Theta_i)$.
The posterior satisfies:
\[
  P(\mathbf{S} \mid \mathbf{D}) \propto
  \prod_{i=1}^n P(\mathbf{D}_i \mid c_i) \cdot P(c_i \mid \Theta_i) \cdot P(\delta_i \mid \Theta_i)
\]
After the update $\Theta_k \to \Theta_k'$, only the $k$-th factor
changes:
\[
  P'(\mathbf{S} \mid \mathbf{D}) \propto
  P(\mathbf{S} \mid \mathbf{D}) \cdot
  \frac{P(c_k \mid \Theta_k') \cdot P(\delta_k \mid \Theta_k')}
       {P(c_k \mid \Theta_k) \cdot P(\delta_k \mid \Theta_k)}
\]

\textit{Step 2: Drift correction.}
For a threshold change $\theta_k \to \theta_k'$, define
$\varepsilon = (\theta_k' - \theta_k)/\theta_k$. The drift posterior
mean updates as $\bar{\mu}_k \leftarrow \bar{\mu}_k - \varepsilon$.
This is a single scalar subtraction: $O(1)$.

\textit{Step 3: Invariance of compliance parameters.}
The compliance rate $c_k$ is independent of the threshold magnitude
(it measures the \emph{proportion} of applicable entities that comply,
not the threshold itself). Therefore $\bar{\alpha}_k$ and
$\bar{\beta}_k$ are invariant under $\mathcal{U}_k$. Similarly,
$\bar{\sigma}_k$ depends on the number of observations and prior
variance, neither of which changes under $\mathcal{U}_k$.

\textit{Step 4: Normalization.}
Under conjugate families, the posterior normalization constant
$B(\bar{\alpha}_k, \bar{\beta}_k)$ is computed analytically and does
not require re-processing the data. The updated posterior is a
properly normalized Beta distribution.

\textit{Total cost:} one scalar subtraction ($\bar{\mu}_k$), zero
data access. $\mathcal{C}^{\mathrm{RSI}}(\mathcal{U}_k) = O(1)$
with respect to $|\mathbf{D}|$ and $n$. \qed
\end{proof}

\subsection{Detailed Proof of Theorem~\ref{thm:bvm} (BvM Consistency)}

\begin{proof}
We prove the result for a single rule $r_i$ (the extension to
$\mathbf{S}$ follows by factorization).

Let $s_1, s_2, \ldots, s_m$ be i.i.d.\ Bernoulli($c_i^*$) compliance
signals, where $c_i^* \in (0,1)$ is the true population compliance rate.

\textit{Step 1: Posterior computation.}
The conjugate posterior is
$Q^*(c_i) = \operatorname{Beta}(\bar{\alpha}_i, \bar{\beta}_i)$ with
$\bar{\alpha}_i = \alpha_i + \sum_{j=1}^m s_j$ and
$\bar{\beta}_i = \beta_i + \sum_{j=1}^m (1-s_j)$.

\textit{Step 2: Posterior mean convergence.}
Define $\bar{s} = m^{-1}\sum_j s_j$. Then:
\[
  \mathbb{E}_Q[c_i] = \frac{\bar{\alpha}_i}{\bar{\alpha}_i + \bar{\beta}_i}
  = \frac{\alpha_i + m\bar{s}}{\alpha_i + \beta_i + m}
  \to \bar{s} \to c_i^* \quad (m \to \infty)
\]
where the first convergence follows from $m \gg \alpha_i + \beta_i$
and the second from the strong law of large numbers.

\textit{Step 3: Posterior variance.}
\[
  \mathrm{Var}_Q[c_i] = \frac{\bar{\alpha}_i \bar{\beta}_i}
  {(\bar{\alpha}_i + \bar{\beta}_i)^2 (\bar{\alpha}_i + \bar{\beta}_i + 1)}
  \sim \frac{c_i^*(1-c_i^*)}{m} \quad (m \to \infty)
\]
This is the Bernoulli variance divided by $m$, confirming
$\sigma[c_i] \sim 1/\sqrt{m}$.

\textit{Step 4: Asymptotic normality.}
By the central limit theorem, $\sqrt{m}(\bar{s} - c_i^*)
\xrightarrow{d} \mathcal{N}(0, c_i^*(1-c_i^*))$. Since
$\mathbb{E}_Q[c_i] \to \bar{s}$ and $\mathrm{Var}_Q[c_i] \to
c_i^*(1-c_i^*)/m$, we obtain:
\[
  \sqrt{m}(\mathbb{E}_Q[c_i] - c_i^*) \xrightarrow{d}
  \mathcal{N}(0, \mathcal{I}(c_i^*)^{-1})
\]
where $\mathcal{I}(c_i^*) = [c_i^*(1-c_i^*)]^{-1}$ is the Fisher
information of the Bernoulli model. This is the classical
Bernstein-von Mises theorem applied to the Beta-Bernoulli conjugate
pair \citep{van2000asymptotic}. \qed
\end{proof}

%% ─────────────────────────────────────────────────────────────
%% A.6 — SIGNAL DISCRETIZATION BIAS
%% ─────────────────────────────────────────────────────────────

\subsection{Signal Discretization and Calibration Bias}

The compliance signal $s_{ji}^{\mathrm{raw}} \in [0,1]$ is unbiased
by construction:
$\mathbb{E}[s_{ji}^{\mathrm{raw}}] = c_{ji}$.
Discretization at threshold $\tau = 0.5$ introduces a systematic bias.
For an entity with true compliance $c_{ji}$, the probability that the
discretized signal equals 1 is:
\[
  P(s_{ji}^{\mathrm{disc}} = 1) =
  1 - I_{0.5}(c_{ji} \cdot \kappa,\, (1{-}c_{ji}) \cdot \kappa)
\]
where $I_x(a,b)$ is the regularized incomplete Beta function. When
$c_{ji} > 0.5$, this probability exceeds $c_{ji}$ (the discrete signal
overestimates compliance). When $c_{ji} < 0.5$, it underestimates.
The bias is largest for $c_{ji}$ far from 0.5 and decreases as
$\kappa \to \infty$.

This explains the calibration misses in Table~\ref{tab:calibration}:
rules with true population compliance far from 0.5 (R3\_IMF at 0.876,
R7\_DECL at 0.666) exhibit the largest bias, while rules near 0.5
(R4\_TPU at 0.345, R8\_BANK at 0.491) are well-calibrated.

A natural mitigation is to use continuous signals directly with a
Beta-likelihood formulation, eliminating discretization entirely while
preserving conjugacy under moment-matching approximations. This
extension is developed in future work.

%% ─────────────────────────────────────────────────────────────
%% A.7 — PRIOR SENSITIVITY ANALYSIS
%% ─────────────────────────────────────────────────────────────

\subsection{Prior Sensitivity Analysis}

To assess robustness to prior misspecification, we evaluate RSI under
five prior configurations ranging from domain-informed to deliberately
adversarial:

\begin{center}
\renewcommand{\arraystretch}{1.15}
\begin{tabular}{llcc}
\toprule
\textbf{Configuration} & \textbf{Prior} & $\mathbb{E}[c_i]$ & \textbf{Rationale} \\
\midrule
Correct       & $\mathrm{Beta}(7, 3)$ & 0.70 & Domain-informed (OTR) \\
Uniform       & $\mathrm{Beta}(1, 1)$ & 0.50 & Maximum ignorance \\
Inverted      & $\mathrm{Beta}(3, 7)$ & 0.30 & Deliberately wrong direction \\
Strong wrong  & $\mathrm{Beta}(20, 2)$ & 0.91 & Strong confidence, wrong value \\
Weak          & $\mathrm{Beta}(2, 2)$ & 0.50 & Weak, symmetric \\
\bottomrule
\end{tabular}
\end{center}

Entity-level F1 is perfectly invariant (0.741 across all
configurations) because entity scoring is deterministic:
$\mathrm{NC}_{ji} = 1 - s_{ji}^{\mathrm{raw}}$ does not depend on
the prior. Population posteriors $\mathbb{E}[c_i]$ vary by at most
0.027 across configurations, confirming that with $n > 400$
observations per rule, the likelihood dominates the prior. This
validates the prior robustness corollary of Theorem~\ref{thm:bvm}:
different initial beliefs converge to the same assessment as evidence
accumulates.